% ============================================================================
% Title:        Introduction (Literature review)
% Author:       Nathaniel Thomas
% Affiliation:  Texas A&M University, Department of Nuclear Engineering
% Email:        nathaniel@tamu.edu
% Date Created: 3/11/2026
% Version:      1.0
%
% Description:
%   A .tex source file containing the transcribed introduction to Kyle and
%   Steven Raiman's literature review on flouride salt purification.
%
% =============================================================================

\section{History of Molten Fluoride Salt Purification}
\label{sec: rev history}

The concept of using molten fluoride salt loops has existed since at
least the late 1940s, and the first experiments conducted toward
building one were performed as part of the Aircraft Reactor
Experiment (ARE) at Oak Ridge National Lab starting in 1949,
practically ending in 1957, and officially concluding in 1961.
In 1952, ARE workers decided that the fluoride compounds purchased to
create experimental molten fluoride salt eutectics were of an
unacceptable level of purity, as oxide impurities in the salt were
causing a sludge to form as experiments proceeded which clogged pipes
\cite{rev1}. An initial proposed purification process consisted of heating
the salt mix to \SI{600}{\celsius} under vacuum to remove any water from the
system, followed by a hydrogen gas (\ch{H2}) sparge at temperature to
remove excess fluorine from the system \cite{rev1}. Later that same year,
they discovered that a lab scale process of treating the salt with an
\ch{HF-H2} atmosphere could be scaled up and applied to reduce the level
of sulfur contaminants in the salt to acceptable \cite{rev2}. This process
would later be refined extended to also remove most oxide
contaminants present in the salt, creating the method of purification
through hydrofluorination that will be described in more detail in
\Cref{sec: hf sparging}. Later that same year, it was discovered that adding
reducing agents to the molten salt, long thought to only have
deleterious effects, could actually help prevent corrosion, so long
as an excessive amount was not added \cite{rev3}. It was also discovered that
a longer purification process and an additional step involving
sparging with pure hydrogen gas could be used to further reduce
corrosion \cite{rev3}. In 1953 it was found that the hydrogen sparge was
reducing corrosion by reducing metals dissolved in the mix during
fluorination to their metallic state and removing excess fluorine
from the system \cite{rev4}. In 1954 it was found that of these metals,
nickel would reduce out quickly by hydrogen sparging while iron and
chromium would be rate limiting and that the presence of these metals
in in a molten salt solution could cause UF4 from fuel to be reduced
to \ch{UF3} \cite{rev5}. This is also around the time they started experimenting
with electrolytic means of purifying the salt \cite{rev5}. In 1955 it was
discovered that hydrogen sparging alone would not reduce out
appreciable concentrations of chromium at a practical rate and that
reducing agents would need to be added \cite{rev6}. A new type of reaction
vessel utilizing copper lined stainless steel in place of nickel was
also put into use in order to scale up salt production \cite{rev6}.
The lessons learned from the initial years of the ARE were then taken
into a new project known as the Molten Salt Reactor Experiment (MSRE)
with a focus shifted towards conducting experiments directed at
making a reactor for producing commercial power. This project,
beginning in from 1957, starting operation fueled in 1965, concluding
fueled operation in 1969, and continuing experiments until 1975,
provided the largest single source of data for the performance,
requirements, and behaviors of a commercial power plant scale molten salt loop.
Beginning in 1957, it was noticed that the new salt chosen (\ch{Li2BeF4})
was substantially more difficult to purify than the old salt
(\ch{NaZrF5}), requiring longer sparge times \cite{rev7},
\cite{rev8}. In 1959, tests
using molten ammonium bifluoride (\ch{NH4HF2}) showed that it performs an
adequate job of converting metal oxides into fluorides, making it a
viable purification method \cite{rev9}. In 1963 a standard fuel, flush, and
cooling salt purification operation procedure had been established,
once again using the proven method of hydrofluorination in nickel
lined vessels and averaged a processing time of \SI{164}{\hour} per batch
\cite{rev10}. This operation remained essentially unchanged throughout the
period of fueled operation \cite{rev10}, \cite{rev11}.
During and after the conclusion of fueled operation experiments into
purifying salt eutectics shifted from removing impurities to removing
fission products produced during operation. Removal of fission
products was largely accomplished through fluorination. In 1972,
experiments were underway to test the viability of a “frozen-wall
fluorinator” \cite{rev12}, \cite{rev13}. Studies into the viability
of this design
had been underway throughout the operating and post operating period
\cite{rev12}, \cite{rev14}.  The purpose of this is rather
self-explanatory in that a
thin layer of fluoride salt would cover walls of the fluorination
vessel and kept below its melting point to remain solid, thus
creating a protective salt layer and allowing for more aggressive
fluorination agents (notably fluorine gas or F2) to be used with
minimal attack on structural metals \cite{rev12}, \cite{rev13}. The
complication of
this design is that it requires internal heating, for which
autoresistance heating was tested, but radiofrequency induction
heating was also considered \cite{rev12}, \cite{rev13}, \cite{rev15}.
Although designed for
the aggressive removal of fission products from used salt, this
design, if proven viable, could be usable in purifying unused salt as
well. Also tested was the distillation of used salt to remove fission
products; however, this proved inefficient \cite{rev11}. After this report
little innovation into the salt purification process was found. This
is likely due to the project being essentially dead by this point, so
funding and manpower were minimal until all work on the project ceased in 1975.

\section{Types of Molten Salt and their Properties}
\label{sec: salt properties}

\section{Effects of Contaminants of Molten Salt Performance}
\label{sec: contaminants}

\subsection{Oxygen}
\label{sec: oxygen}

\subsection{Sulfur}
\label{sec: sulfur}

\subsection{Structural Metals}
\label{sec: structural metals}

\section{Purification Methods}
\label{sec: purification methods}

\subsection{\ch{HF}-\ch{H2} Sparging}
\label{sec: hf sparging}

\subsubsection{Chemistry}
\label{sec: hf chemistry}

\subsubsection{Kinetics}
\label{sec: hf kinetics}

\subsubsection{Reagent Properties and Hazards}
\label{sec: hf reagent properties}

\subsubsection{Reagent Costs}
\label{sec: hf reagent costs}

\subsection{\ch{NF3}-\ch{H2} Sparging/\ch{NF3} Baking}
\label{sec: nf3 sparging}

\subsubsection{Chemistry and Thermodynamics}
\label{sec: nf3 chemistry and thermodynamics}

\subsubsection{Kinetics}
\label{sec: kinetics}

\subsubsection{Reagent Properties and Hazard}
\label{sec: nf3 hazards}

\subsubsection{Reagent Costs}
\label{sec: nf3 reagent costs}

\subsection{\ch{NF4HF2} Baking}
\label{sec: nf4hf2 baking}

\subsubsection{Chemistry}
\label{sec: nf4hf2 chemistry}

\subsubsection{Kinetics}
\label{sec: nf4hf2 kinetics}

\subsubsection{Reagent Properties and Hazards}
\label{sec: nf4hf2 hazards}

\subsection{Electrolytic Purification}
\label{sec: electrolytic purification}

\section{Additional Steps}
\label{sec: additional steps}

\subsection{Continued H addition}
\label{sec: continued h addition}

\subsection{Reducing metal additions}
\label{sec: reducing metal additions}

\pagebreak

\subsection{Filtration}
\label{sec: filtration}

\subsection{Corrosion in Structural Materials}
\label{sec: corrosion}

\subsection{Corrosion Due to Hydrofluorinating Reagents}
\label{sec: corrosion due to hf}

\subsection{Embrittlement Due to Hydrogen}
\label{sec: embrittlement}

\subsection{Corrosion Due to Molten Salt}
\label{sec: corrosion due to salt}

\section{Compatibility of Candidate Structural Materials}
\label{sec: compatability of candidate}

\subsection{Ni Alloys}
\label{sec: ni alloys}

\subsection{Steels}
\label{sec: steels}

\subsection{Copper}
\label{sec: copper}

\subsection{Carbon}
\label{sec: carbon}

\section{References}
\label{sec: references}

\printbibliography[heading=subbibliography]
\end{refsection}
